% =====================================================================================
%  APPENDIX: the full statistical re-analysis promised by Sec. 6.5.  WRITTEN 2026-09-18.
%  app:stats was referenced by sec_6.tex and defined nowhere.
%  Every number below is read from release/data/stats_resource.json (stats_resource.py,
%  2026-09-18T18:53Z) and was re-read from that file while this appendix was written.
%  NOTE: the Williams statistic is t = 0.9449, i.e. 0.94 -- Sec. 6.5 previously printed
%  0.95 and Fig. 2's caption printed 0.94; the file says 0.94 and both now agree with it.
% =====================================================================================

\section{The resource statistics, stratified}
\label{app:stats}

Section~\ref{sec:nogo:chi} reports the correlation evidence within the design that produced it. This
appendix gives the design, the per-stratum coefficients and the three tests, so that the retraction of
the two pooled coefficients of versions~1--2 can be checked rather than taken.

\paragraph{The lattice design.} The $30$ Hubbard points are six $(L,N)$ strata swept over the same five
on-site repulsions $U=1,2,4,8,16$; they are not $30$ independent observations, and the effective number
of independent sweeps is six. Table~\ref{tab:strata} gives each stratum. The rank correlation of
$\mathcal F_1$ with the static support $|\mathcal S|_{\varepsilon_w}$ is $-1$ in all six, exactly; that
of $\chi$ runs from $+0.894$ to $+1.000$ with mean $+0.948$, its departures from unity being ties in
$\chi$ and not inversions.

\begin{table*}[tp]\centering\small
\caption{\textbf{The six Hubbard strata.} Each is a five-point sweep in $U$ at fixed $(L,N)$. The
one-sided exact permutation probability of a perfect reversal within one stratum of five distinct
values is $1/120=0.0083$; ties in $\chi$ raise it. ``distinct $\chi$'' counts how many of the five
$\chi$ values differ, which is what caps $\rho(\chi,|\mathcal S|)$ below unity in four of the
six.}\label{tab:strata}
\begin{tabular}{cccccc}
\toprule
$L$ & $N$ & $\rho(\mathcal F_1,|\mathcal S|)$ & $\rho(\chi,|\mathcal S|)$ & $\chi$ over $U=1,2,4,8,16$ & distinct $\chi$\\
\midrule
$6$  & $6$  & $-1.000$ & $+1.000$ & $12,10,6,5,4$   & $5$\\
$6$  & $4$  & $-1.000$ & $+0.894$ & $11,10,9,9,9$   & $3$\\
$8$  & $8$  & $-1.000$ & $+0.949$ & $11,11,8,8,7$   & $3$\\
$8$  & $6$  & $-1.000$ & $+0.975$ & $16,16,15,14,13$& $4$\\
$10$ & $10$ & $-1.000$ & $+0.975$ & $17,14,10,6,6$  & $4$\\
$10$ & $8$  & $-1.000$ & $+0.894$ & $13,13,13,11,9$ & $3$\\
\bottomrule
\end{tabular}
\end{table*}

\paragraph{Three tests, and what each assumes.} The stratified Kendall statistics are $\tau=-1.000$ for
$\mathcal F_1$, with $0$ concordant and $60$ discordant within-stratum pairs and no ties, and
$\tau=+1.000$ for $\chi$, with $50$ concordant, $0$ discordant and $10$ tied. Because each
$\mathcal F_1$ stratum is a perfect reversal of five distinct points, the exact one-sided stratified
permutation probability is the product of the six per-stratum chances,
$(1/120)^{6}=3.35\times10^{-13}$, over a relabelling space of $2.99\times10^{12}$; the same test
applied to $\chi$, which has ties, gives $1.93\times10^{-10}$. Fisher's combination of the six
per-stratum probabilities is conservative here and returns $6.56\times10^{-8}$ and
$1.14\times10^{-5}$ ($\chi^{2}=57.4$ and $44.7$ on $12$ degrees of freedom); we quote the permutation
values and record Fisher's beside them. All are \emph{one-sided}---a perfect reversal is one of the
$120$ orderings, not two---and all take a product across strata, so all assume the six strata
independent. They are not six independent experiments: they are three lengths crossed with two
fillings, on one code path and one $U$ grid. What the numbers establish is that the direction is
reproduced six times out of six, not that it generalizes beyond this family.

\paragraph{Why the pooled coefficients are withdrawn, both of them.} Pooled over the same $30$ points
the coefficients are $\rho(\chi,|\mathcal S|)=+0.569$ and
$\rho(\mathcal F_1,|\mathcal S|)=-0.109$, the second with a nominal two-sided $p=0.57$ and the first
with $p=0.0010$. Pooling six deterministic five-point sweeps with different offsets destroys the
within-stratum structure that carries the signal; it is a Simpson reversal in both cases. It happens to
destroy the signal in our favour on $\chi$ and against us on $\mathcal F_1$, and that is exactly why
neither may be kept: a rule that is applied to one coefficient and not to the other is not a rule.

\paragraph{The molecular set.} The $n=38$ molecular points are $19$ molecules at two geometries, i.e.\
dependent pairs. The intraclass correlation of $\log_2|\mathcal S|$ within a molecule is $0.214$, a
design effect of $1.214$ and a design-corrected effective $n$ of $31.3$. Ties are pervasive: $29$ of
the $38$ ($76\%$) lie in some tied group of $|\mathcal S|_{\varepsilon_w}$ and $11$ are pinned at the
floor $|\mathcal S|=1$, which caps the attainable Spearman coefficient at $\rho_{\max}=0.9863$;
against that ceiling $\chi$ reaches $91.4\%$ and $\mathcal F_1$ $87.4\%$. Under a molecule-level
cluster permutation ($10^{6}$ resamples) both coefficients are beyond the resolution floor of the
test, $p\le1.0\times10^{-6}$. The comparison versions~1--2 rested on is not resolvable, and three
procedures agree that it is not: $\rho=0.9011$ for $\chi$ against $0.8618$ for $\mathcal F_1$ differ by
$0.0393$; a molecule-level cluster bootstrap ($2\times10^{4}$ resamples) puts a $95\%$ interval on the
difference of $[-0.070,+0.189]$, which contains zero; Williams's test for dependent correlations gives
$t=0.94$ on $35$ degrees of freedom, $p=0.35$, and Steiger's $z=0.94$, $p=0.35$. Both parametric tests
assume bivariate normality and untied ranks while $76\%$ of these points sit in a tie, which is why the
bootstrap is quoted first; at the cluster-level effective $n$ of $19$, Williams gives $p=0.53$. Pooled
over all $74$ rows the pair is $0.722$ against $0.599$ and is likewise unresolved. Recomputed over the
$71$ \emph{distinct} states---three chain rows duplicate three Hubbard rows---it is $0.746$ against
$0.592$.

\paragraph{One artefact of the deposit, recorded because it changes a digit.} The pooled
$\rho(\mathcal F_1,|\mathcal S|)$ stored in the deposit is $0.5991984$ and recomputing it from the
per-species values gives $0.5992028$. The difference is not a bug in either: it is a tie in
$\mathcal F_1$ created when two rows that are the same state are stored at different precisions. The
quantity is quoted to three figures throughout for that reason.
